% answers_short.tex —— 2011 年·陕西卷·文科 · 简短答案速查
% 由 tools/make_short_answers.py 自 answers.tex 的 \gksolhead/\gkanswer 提取，
% 并与 2011/2011_数学真题_陕西文_答案要点.md 逐题比对后生成，勿手改（改源 MD 后重跑）。

\begin{gknotice}
\textbf{使用说明}：本节只给最终答案便于快速核对，完整推导见解析版。答案为编写者独立推导，\textbf{非官方标准答案}。
\end{gknotice}

\gkpart{选择题}
\begin{center}\begin{tabular}{ccccccccccc}
\toprule
题号 & 1 & 2 & 3 & 4 & 5 & 6 & 7 & 8 & 9 & 10 \\
\midrule
答案 & D & C & B & B & A & C & B & C & A & D \\
\bottomrule
\end{tabular}\end{center}

\gkpart{填空题}
\begin{itemize}[leftmargin=2.6em]
  \item[\textbf{11.}] $-2$。
  \item[\textbf{12.}] $1$。
  \item[\textbf{13.}] $5+6+\cdots+13=81$。
  \item[\textbf{14.}] $3$ 或 $4$。
  \item[\textbf{15.}] A $a\leqslant3$；B $2$；C $1$。
\end{itemize}

\gkpart{解答题}
\begin{itemize}[leftmargin=2.6em]
  \item[\textbf{16.}] (1) 证明见解析；(2) $\dfrac52$
  \item[\textbf{17.}] (1) $\dfrac{x^2}{25}+\dfrac{y^2}{16}=1$；(2) $\left(\dfrac32,-\dfrac65\right)$
  \item[\textbf{18.}] 余弦定理：三角形任何一边的平方等于其他两边平方的和减去这两边与它们夹角的余弦的积的两倍；证明见解析。
  \item[\textbf{19.}] (1) $x_k=x_{k-1}-1$（$2\leqslant k\leqslant n$）；(2) $\dfrac{\mathrm e\left(1-\mathrm e^{-n}\right)}{\mathrm e-1}$
  \item[\textbf{20.}] (1) $0.44$；(2) $L_1$ 各时间段频率依次为 $0.1$，$0.2$，$0.3$，$0.2$，$0.2$；$L_2$ 依次为 $0$，$0.1$，$0.4$，$0.4$，$0.1$；(3) 甲应选择路径 $L_1$，乙应选择路径 $L_2$
  \item[\textbf{21.}] (1) $g(x)$ 在 $(0,1)$ 上单调递减，在 $(1,+\infty)$ 上单调递增，最小值为 $g(1)=1$；(2) 当 $0<x<1$ 时 $g(x)>g\left(\dfrac1x\right)$；当 $x=1$ 时两者相等；当 $x>1$ 时 $g(x)<g\left(\dfrac1x\right)$；(3) $0<a<\mathrm e$
\end{itemize}
